\documentclass[11pt]{article}
\usepackage[margin=1in]{geometry}
\usepackage{amsmath,amssymb,amsthm}
\usepackage{hyperref}

\newcommand{\Vhat}{\hat V}
\newcommand{\Qhat}{\hat Q}
\newcommand{\pihat}{\hat\pi}
\newcommand{\PY}{P_{Y,t}}
\newcommand{\E}{\mathbb{E}}

\theoremstyle{definition}
\newtheorem{definition}{Definition}

\title{Why the Extended Theorem 3 lower bound\\ collapses to the trivial $V^{\pihat^\star}-V^\star\ge 0$}
\date{\today}

\begin{document}
\maketitle

\begin{abstract}
Extended Theorem~3 is two-sided: it states both an upper bound (Eq.~21) and a
symmetric lower bound (Eq.~27) on the sub-optimality gap
$V^{\pihat^\star}_t-V^\star_t$. The upper bound is informative and computable
without the true optimal policy. The lower bound, however, when instantiated
\emph{within the framework} (i.e.\ without re-introducing the true model's optimal
value $V^\star$ or policy $\pi^\star$), can only certify the trivial statement
$V^{\pihat^\star}_t-V^\star_t\ge 0$. This note explains \emph{why}, with the key
asymmetry isolated in one place.
\end{abstract}

\section{Notation and the two error streams}
We compare a \emph{true} model with optimal value/policy $(V^\star,\pi^\star)$ and
an \emph{approximate} (AIS) model with optimal value/policy
$(\Vhat^\star,\pihat^\star)$. The policy actually deployed is the AIS-optimal
policy lifted to the true system, $\pi_t(h_t)=\pihat^\star_t(\sigma_t(h_t))$, whose
true value is $V^{\pihat^\star}_t$. Theorem~3 tracks \emph{two signed error
streams}, both measured against the shared baseline $\Vhat^\star$:
\begin{align}
\alpha^{\pihat^\star}_t(h_t) &= V^{\pihat^\star}_t(h_t)-\Vhat^\star_t(\sigma_t(h_t))
&&\text{(policy-evaluation stream),}\\
\alpha^{\star}_t(h_t) &= V^{\star}_t(h_t)-\Vhat^\star_t(\sigma_t(h_t))
&&\text{(optimality stream).}
\end{align}
Subtracting the two cancels the baseline and returns exactly the quantity we care
about,
\begin{equation}
\alpha^{\pihat^\star}_t(h_t)-\alpha^{\star}_t(h_t)
= V^{\pihat^\star}_t(h_t)-V^{\star}_t(h_t),
\label{eq:bridge}
\end{equation}
the ``bridge identity.'' Each stream has a state--action ($Q$-like) companion,
\begin{equation}
\beta^{\pihat^\star}_t(h_t,a_t)=Q^{\pihat^\star}_t(h_t,a_t)-\Qhat^\star_t(\sigma_t(h_t),a_t),
\qquad
\beta^{\star}_t(h_t,a_t)=Q^{\star}_t(h_t,a_t)-\Qhat^\star_t(\sigma_t(h_t),a_t).
\end{equation}
The single one-step mismatch shared by both streams is
\begin{equation}
\Delta_t(h_t,a_t)=
\Bigl[c_t(h_t,a_t)+\textstyle\int \Vhat^\star_{t+1}\,dP^\sigma_t\Bigr]
-\Bigl[\hat c_t(\sigma_t(h_t),a_t)+\textstyle\int \Vhat^\star_{t+1}\,d\hat P^\sigma_t\Bigr],
\label{eq:delta}
\end{equation}
i.e.\ the cost-plus-kernel mismatch evaluated on the \emph{same} baseline
$\Vhat^\star_{t+1}$, propagated by the two models' kernels. Crucially, both streams
use the \emph{same} $\Delta_t$: because $\pihat^\star$ is AIS-optimal we have
$\Vhat^{\pihat^\star}=\Vhat^\star$, so the policy-evaluation mismatch coincides
with the optimality mismatch.

\section{Sub-solutions and super-solutions}
The exact streams obey one-step recursions (Lemmas~1--2). For example the
\emph{exact} optimality stream $\tilde\alpha^\star_t=V^\star_t-\Vhat^\star_t$
satisfies
\begin{equation}
\tilde\beta^\star_t(h_t,a_t)=\Delta_t(h_t,a_t)+\int \tilde\alpha^\star_{t+1}\,d\PY ,
\label{eq:exactrec}
\end{equation}
and an analogous equality holds for the policy-evaluation stream. A
\emph{certificate} replaces these equalities with one-sided inequalities. Two
directions are possible.

\begin{definition}[Super-solution]
A function is a \emph{super-solution} if it sits \textbf{above} the exact stream,
$\alpha_t\ge\tilde\alpha_t$, and is produced by turning the exact recursion into a
``$\ge$'': $\beta_t\ge\Delta_t+\int\alpha_{t+1}\,d\PY$. A super-solution
\emph{over-estimates} the error.
\end{definition}

\begin{definition}[Sub-solution]
A function is a \emph{sub-solution} if it sits \textbf{below} the exact stream,
$\alpha_t\le\tilde\alpha_t$, and is produced by turning the exact recursion into a
``$\le$'': $\beta_t\le\Delta_t+\int\alpha_{t+1}\,d\PY$. A sub-solution
\emph{under-estimates} the error.
\end{definition}

\paragraph{Why the directions matter.} From the bridge identity
\eqref{eq:bridge}, $V^{\pihat^\star}_t-V^\star_t=\tilde\alpha^{\pihat^\star}_t-\tilde\alpha^\star_t$.
To get a \emph{valid} bound we replace each exact stream by a certificate that
pushes the difference the right way:
\begin{itemize}
\item \textbf{Upper bound:} make the first term larger and the second smaller, i.e.\
use a super-solution for $\alpha^{\pihat^\star}$ and a sub-solution for
$\alpha^\star$. Then
$V^{\pihat^\star}_t-V^\star_t=\tilde\alpha^{\pihat^\star}_t-\tilde\alpha^\star_t
\le \alpha^{\pihat^\star}_t-\alpha^\star_t$.
\item \textbf{Lower bound:} the reverse---use a sub-solution for
$\alpha^{\pihat^\star}$ and a super-solution for $\alpha^\star$. Then
$V^{\pihat^\star}_t-V^\star_t\ge \alpha^{\pihat^\star}_t-\alpha^\star_t$.
\end{itemize}

\paragraph{The anchoring action.} A recursion is only closed once we say at
\emph{which action} $\beta_t(h_t,\cdot)$ is evaluated to form $\alpha_t(h_t)$. This
``anchor'' is where all the freedom---and all the trouble---lives.

\section{The lower-bound hypotheses of Theorem 3}
For the lower bound, Theorem~3 asks for candidate functions satisfying, for all
$t$,
\begin{align}
\beta^{\pihat^\star}_T \le \Delta_T \le \beta^\star_T, &&\text{(terminal, Eq.~22)}\label{eq:T}\\
\beta^{\pihat^\star}_t(h_t,a_t) \le \Delta_t(h_t,a_t)+\int \alpha^{\pihat^\star}_{t+1}\,d\PY,
&&\text{(sub-solution recursion, Eq.~23)}\label{eq:23}\\
\beta^{\star}_t(h_t,a_t) \ge \Delta_t(h_t,a_t)+\int \alpha^{\star}_{t+1}\,d\PY,
&&\text{(super-solution recursion, Eq.~24)}\label{eq:24}\\
\alpha^{\pihat^\star}_t(h_t) \le \beta^{\pihat^\star}_t\bigl(h_t,\pihat^\star_t(\sigma_t(h_t))\bigr),
&&\text{(anchor at $\pihat^\star$, Eq.~25)}\label{eq:25}\\
\alpha^{\star}_t(h_t) \ge \beta^{\star}_t\bigl(h_t,\pihat^\star_t(\sigma_t(h_t))\bigr).
&&\text{(anchor at $\pihat^\star$, Eq.~26)}\label{eq:26}
\end{align}
The conclusion is the lower bound
\begin{equation}
V^{\pihat^\star}_t(h_t)-V^\star_t(h_t)\ \ge\ \alpha^{\pihat^\star}_t(h_t)-\alpha^{\star}_t(h_t).
\tag{27}\label{eq:27}
\end{equation}

\section{The collapse}
We want the \emph{tightest} (largest, hence most informative) admissible value of
$\alpha^{\pihat^\star}_t-\alpha^\star_t$. Maximize each piece subject to its
constraints.

\paragraph{Largest admissible $\alpha^{\pihat^\star}_t$.} Chaining \eqref{eq:25}
then \eqref{eq:23} (both at the anchor $\pihat^\star$) gives the upper envelope
\begin{equation}
\alpha^{\pihat^\star}_t(h_t)\ \le\ \Delta_t\bigl(h_t,\pihat^\star\bigr)+\int \alpha^{\pihat^\star}_{t+1}\,d\PY,
\label{eq:rec_ph}
\end{equation}
so the largest admissible choice attains equality in \eqref{eq:rec_ph}.

\paragraph{Smallest admissible $\alpha^{\star}_t$.} Chaining \eqref{eq:26} then
\eqref{eq:24} (again at the anchor $\pihat^\star$) gives the lower envelope
\begin{equation}
\alpha^{\star}_t(h_t)\ \ge\ \Delta_t\bigl(h_t,\pihat^\star\bigr)+\int \alpha^{\star}_{t+1}\,d\PY,
\label{eq:rec_star}
\end{equation}
so the smallest admissible choice attains equality in \eqref{eq:rec_star}.

\paragraph{They are the same recursion.} Compare \eqref{eq:rec_ph} and
\eqref{eq:rec_star} at equality. They have
\begin{itemize}
\item the \textbf{same anchor action} $\pihat^\star_t(\sigma_t(h_t))$ (forced by
\eqref{eq:25} and \eqref{eq:26});
\item the \textbf{same mismatch} $\Delta_t$ (one shared object, \eqref{eq:delta});
\item the \textbf{same kernel} $\PY$ in the propagation;
\item the \textbf{same terminal value}: \eqref{eq:T} squeezes both
$\beta^{\pihat^\star}_T$ and $\beta^\star_T$ to $\Delta_T$, hence both
$\alpha^{\pihat^\star}_T=\alpha^\star_T=\Delta_T$.
\end{itemize}
A backward induction now gives, for every $t$,
\begin{equation}
\alpha^{\pihat^\star}_t(h_t)\ \equiv\ \alpha^{\star}_t(h_t)
\qquad\Longrightarrow\qquad
\alpha^{\pihat^\star}_t-\alpha^{\star}_t = 0 .
\end{equation}
\emph{Induction.} At $t=T$ both equal $\Delta_T$. If they coincide at $t+1$, then
the integrals $\int\alpha^{\pihat^\star}_{t+1}\,d\PY$ and
$\int\alpha^{\star}_{t+1}\,d\PY$ are equal, and since the remaining terms
($\Delta_t$ at the same anchor) are identical, \eqref{eq:rec_ph} and
\eqref{eq:rec_star} return the same value at $t$. Hence the best lower bound
Theorem~3 yields in-framework is the trivial
\[
V^{\pihat^\star}_t-V^\star_t\ \ge\ 0 .
\]

\section{Why---in one sentence}
The optimality stream $\alpha^\star=V^\star-\Vhat^\star$ involves the \emph{true}
optimal value $V^\star$, which the framework never computes. Its only handle on the
true system is the single policy it knows, $\pihat^\star$. Equation~\eqref{eq:26}
forces the optimality super-solution to be anchored at $\pihat^\star$---but
evaluating the true model at $\pihat^\star$ yields $V^{\pihat^\star}$, \emph{not}
$V^\star$. That is precisely the policy-evaluation stream. So the smallest
super-solution the theorem permits is dragged all the way up to
$\alpha^{\pihat^\star}$; it overshoots the real $\tilde\alpha^\star$ by exactly the
gap, and in the difference the two streams cancel.

\section{Contrast: why the upper bound does \emph{not} collapse}
The upper bound changes \textbf{one thing}: the anchor of the optimality stream.
Its hypothesis (Eq.~20) is
\begin{equation}
\alpha^{\star}_t(h_t)\ \le\ \beta^{\star}_t\bigl(h_t,\pi^\star_t(h_t)\bigr),
\tag{20}
\end{equation}
anchored at the \emph{true} optimal action $\pi^\star$, not $\pihat^\star$. Because
this is a \emph{sub}-solution requirement ($\le$), it is realizable by a
\textbf{minimization}
\begin{equation}
\alpha^{\star}_t(h_t)=\min_{a\in\mathcal C(h_t)}\Bigl[\Delta_t(h_t,a)+\int \alpha^{\star}_{t+1}\,d\PY\Bigr],
\qquad \pi^\star_t(h_t)\in\mathcal C(h_t),
\end{equation}
over a candidate set $\mathcal C$ guaranteed to \emph{contain} $\pi^\star$. Taking a
min over a set containing $\pi^\star$ automatically gives
$\alpha^\star_t\le\beta^\star_t(h_t,\pi^\star)$ \emph{without ever identifying
$\pi^\star$ exactly} (e.g.\ $\mathcal C=$ all actions, or a small base-stock set).
That minimizing action is generally \emph{different} from---and cheaper
than---$\pihat^\star$, so the optimality stream genuinely separates from the
policy-evaluation stream and
\[
U_t=\alpha^{\pihat^\star}_t-\alpha^{\star}_t\ >\ 0 .
\]

\paragraph{The asymmetry, summarized.}
\begin{center}
\begin{tabular}{lll}
\hline
 & optimality stream role & anchor \\
\hline
Upper bound (Eq.~21) & sub-solution ($\le$, realizable by $\min$) & true $\pi^\star$ (via $\mathcal C\ni\pi^\star$) \\
Lower bound (Eq.~27) & super-solution ($\ge$) & forced to $\pihat^\star$ (Eq.~26) \\
\hline
\end{tabular}
\end{center}
A \emph{sub}-solution can be anchored at the unknown $\pi^\star$ through a min over
a candidate set---computable, and producing a real gap. A \emph{super}-solution
(Eq.~26) is pinned to the \emph{same} anchor $\pihat^\star$ as policy evaluation, so
it carries no information beyond it, and the certified lower bound is $0$.

\paragraph{Consequence.} A \emph{non-trivial} lower bound on the gap requires
information outside Extended Theorem~3---either the true optimal policy $\pi^\star$,
or an out-of-framework advantage term such as
$V^{\pihat^\star}_1-Q^\star_1(s,\pihat^\star(s))$ that explicitly uses the true
$V^\star$. Within Theorem~3 alone, only the upper certificate $U$ is informative.

\section{The distinction at a glance: which policy each computation uses}
Both bounds begin from the same bridge identity~\eqref{eq:bridge},
\[
V^{\pihat^\star}_t-V^\star_t=\tilde\alpha^{\pihat^\star}_t-\tilde\alpha^\star_t,
\]
and both reuse the \emph{identical} policy-evaluation stream
$\alpha^{\pihat^\star}$, anchored at the deployed policy $\pihat^\star$. They differ
in exactly one decision---\textbf{which policy anchors the optimality stream}
$\alpha^\star$---and that single choice decides whether the certificate is
informative or empty.
\begin{center}
\begin{tabular}{lccc}
\hline
computation & optimality-stream anchor & policy it compares $\pihat^\star$ against & value \\
\hline
upper $U_t$ & $\pi^\star$\, (via $\min_{a\in\mathcal C}$, $\mathcal C\ni\pi^\star$) & the \emph{true optimal} $\pi^\star\neq\pihat^\star$ & $>0$ \\
lower $L_t$ & $\pihat^\star$\, (forced by Eq.~\eqref{eq:26}) & $\pihat^\star$ \emph{itself} & $=0$ \\
\hline
\end{tabular}
\end{center}

\paragraph{Upper bound is nonzero.} The optimality sub-solution is evaluated at the
\emph{true optimal} action $\pi^\star$ (reached implicitly by minimizing over a
candidate set that contains it). Since $\pi^\star\neq\pihat^\star$ at the states
that matter, the two streams track value functions of \emph{two different
policies}; their difference is the genuine sub-optimality gap, so
$U_t=\alpha^{\pihat^\star}_t-\alpha^\star_t>0$.

\paragraph{Lower bound is zero.} The optimality super-solution is pinned by
Eq.~\eqref{eq:26} to $\pihat^\star$---the \emph{same} policy the evaluation stream
already uses. With one shared policy, one shared mismatch $\Delta_t$, one shared
kernel, and one shared terminal, the two streams reconstruct the \emph{same}
recursion and cancel pointwise, leaving the empty statement $L_t\equiv 0$.

\paragraph{Numerical confirmation.} In the inventory experiment
($T=100$, $\gamma=0.75$, base-stocks $\sigma^\star=2$, $\hat\sigma^\star=3$) the
companion notebook computes the $\pihat^\star$-anchored optimality super-solution as
an \emph{independent} backward recursion and finds it bit-exactly equal to
$\alpha^{\pihat^\star}$ at every $(t,s)$, so $L(s)\equiv 0$. The upper certificate,
whose optimality stream instead reaches $\pi^\star$, is strictly positive over the
operating region $s\in[-7,3]$: $U_{\max}=2.06$ and $U_{\mathrm{res}}=2.80$ (true
gap $2.01$). The two outcomes differ \emph{only} because of the anchoring policy.

\end{document}
